\documentclass[11pt,leqno]{amsart}
%%PAGE SETUP

\sloppy
\usepackage[utf8]{inputenc}
\usepackage{amssymb}
\usepackage{amsfonts}
\usepackage{amscd}
\usepackage{amstext}
\usepackage{latexsym,amssymb}
\usepackage{amsmath}
\usepackage{graphicx}
\usepackage{color}
\usepackage{palatino}
%\usepackage{showkeys}


\newcommand{\R}{\mathrm{I\!R\!}}
%\newcommand{\C}{\mathbb{C}}
%\newcommand{\Rn}{\RR^{n}}
%\newcommand{\Rd}{\RR^{2}}
\newcommand{\N}{I\hspace{-0.7ex} N}
\newcommand{\cc}{\c{c}}
\newcommand{\ca}{\c{c}\~a}

\usepackage{epsfig}
\newcommand{\co}{\c{c}\~o}
\newcommand{\disp}{\displaystyle}
\newcommand{\caixa}{\rule{1.3ex}{.7em}}
\newcommand{\peq}[1]{\scriptstyle #1}

%%%%%%%%%%%%%%%%
\usepackage[numbers,sort&compress]{natbib}
%%%%%

\usepackage{natbib}
\usepackage{graphicx}
\newtheorem{theorem}{Theorem}[section]
\newtheorem{corollary}{Corollary}
%\newtheorem*{main}{Main Theorem}
\newtheorem{lemma}[theorem]{Lemma}
\newtheorem{proposition}[theorem]{Proposition}
\newtheorem{conjecture}{Conjecture}
\newtheorem{ex}{Example}
%\theoremstyle{definition}
\newtheorem{definition}[theorem]{Definition}
\newtheorem{remark}[theorem]{Remark}
%\newtheorem*{notation}{Notation}

%\newtheorem{lemma}{Lemma}[section]
%\newtheorem{proposition}{Proposition}[section]
%\newtheorem{theorem}{Theorem}[section]
%\newtheorem{remark}{Remark}[section]
%\newtheorem{definition}{Definition}[section]
%\newtheorem{corollary}{Corollary}[section]
%\newtheorem{ex}{Example}[section]

%\pagestyle{empty}


% AMS proof environment with explicit spacing after the proof heading.
\makeatletter
\renewenvironment{proof}[1][\proofname]{%
  \par
  \pushQED{\qed}%
  \normalfont
  \topsep6\p@\@plus6\p@\relax
  \trivlist
  \item[\hskip\labelsep\itshape #1\@addpunct{.}]%
  \hspace{0.5em}\ignorespaces
}{%
  \popQED
  \endtrivlist\@endpefalse
}
\makeatother

\renewcommand{\thefootnote}{\fnsymbol{footnote}}


%%\date{\today}
\title[Small Positive Solutions for Singular Mean Curvature Equations]
{Small Positive Solutions for Singular Prescribed Mean Curvature Equations with Convection}

\author{Anderson L. A. de Araujo}
\address[Anderson L. A. de Araujo]{Departamento de Matem\'atica, Universidade Federal de Vi\c{c}osa, Avenida Peter Henry Rolfs, s/n, 36570-900, Vi\c{c}osa, MG, Brazil}
\email{anderson.araujo@ufv.br}

\author{Kamila F. L. Madalena}
\address[Kamila F. L. Madalena]{Departamento de Matem\'atica, Universidade Federal de Itajub\'a, 35903-087 - Itabira - MG, Brazil}
\email{kamilalobo@unifei.edu.br}

\date{}


% Page layout: controlled margins compatible with amsart.
% The former manual \textwidth/\oddsidemargin/\topmargin adjustments
% were removed because they distort the AMS page layout.
\usepackage[
  letterpaper,
  textwidth=19.1cm,
  top=1.8cm,
  bottom=2.0cm,
  includeheadfoot
]{geometry}

\begin{document}

\begin{abstract}
We study a singular prescribed mean curvature
Dirichlet problem with a convection term,
\[
-\operatorname{div}\left(\frac{\nabla u}{\sqrt{1+|\nabla u|^2}}\right)
=\frac{\lambda}{u^\alpha}+f(x,u,\nabla u)
\quad\hbox{in }\Omega,\qquad
u>0\quad\hbox{in }\Omega,\qquad
u=0\quad\hbox{on }\partial\Omega,
\]
where $0<\alpha<1$ and the continuous perturbation satisfies
\[
0\le f(x,t,\xi)\le k_1|t|^{r_1}+k_2|\xi|^{r_2},
\qquad r_1,r_2>1.
\]
We prove the existence of a positive solution for all sufficiently small
$\lambda>0$. The argument is nonvariational and combines a
sub--supersolution scheme with a uniformly elliptic modification of the
mean curvature operator, an order truncation, and a contraction of the
gradient variable in the convection term. A central point is the
construction of the supersolution at the natural singular scale
$t_\lambda=\lambda^{1/(1+\alpha)}$. This scale also permits the gradient
contraction to be removed a posteriori. In particular, the solutions
obtained satisfy the quantitative estimate
\[
\|u_\lambda\|_{C^1(\overline\Omega)}
\le C\lambda^{1/(1+\alpha)},
\]
and hence converge to zero in $C^1(\overline\Omega)$ as
$\lambda\to0^+$. The method is specifically designed to handle
simultaneously the singular zero-order term, the gradient dependence,
and the non-uniform ellipticity of the prescribed mean curvature
operator.
\end{abstract}

\keywords{prescribed mean curvature equation, singular elliptic problem, convection term, sub--supersolutions, gradient truncation, small positive solutions.}
\subjclass[2020]{Primary 35J93; Secondary 35J75, 35J62, 35B45, 35B09.}

\maketitle


\section{Introduction}

We consider the singular Dirichlet problem
\begin{equation}\label{P1}
    \left\{
    \begin{array}{lcc}
        -\operatorname{div}\left(
        \dfrac{\nabla u}{\sqrt{1+|\nabla u|^2}}\right)
        =\dfrac{\lambda}{u^\alpha}+f(x,u,\nabla u)
        &\textup{in}&\Omega,\\[1ex]
        u>0&\textup{in}&\Omega,\\
        u=0&\textup{on}&\partial\Omega,
    \end{array}
    \right.
\end{equation}
where $\Omega\subset\mathbb R^N$, $N\ge2$, is a bounded domain with smooth
boundary, $\lambda>0$, $0<\alpha<1$, and
$f:\Omega\times\mathbb R\times\mathbb R^N\to\mathbb R$ is continuous
and satisfies
\begin{equation}\label{f}
0\le f(x,t,\xi)\le k_1|t|^{r_1}+k_2|\xi|^{r_2},
\qquad k_1,k_2>0,\qquad r_1,r_2>1.
\end{equation}
The purpose of the paper is to establish a branch of positive solutions
which is quantitatively small as $\lambda\to0^+$, while retaining both
the singular forcing $u^{-\alpha}$ and the genuine convection dependence
on $\nabla u$.

The operator in \eqref{P1} is the non-parametric prescribed mean
curvature operator. Its geometric origin and the classical Dirichlet
theory are connected with minimal and capillary surfaces; see, among
others, \cite{Finn,Giusti,Serrin,Lopez}. Existence and multiplicity for
nonlinear prescribed mean curvature equations have been studied through
several methods, including barrier arguments, topological and
variational techniques, and finite-dimensional approximations; see
\cite{CoffmanZ,HabetsO,ObersnelO,Lorca1,Le,CGMS}. Because the operator is
not uniformly elliptic, a priori height and gradient bounds play a central
role in this theory; see, for instance,
\cite{Ripoll,Tersenov,Lorca2}. More recently,
Figueiredo and R\u{a}dulescu \cite{FR} treated positive solutions under
critical exponential growth, illustrating that the mean curvature
operator can also be handled in borderline nonlinear regimes. Recent
work has further developed the PDE and BV theory of non-parametric mean
curvature equations with rough data and nonlinear zero-order terms.
Giachetti, Oliva and Petitta \cite{GOP2024}, for instance, obtained
bounded solutions for nonlinear mean curvature problems allowing a
factor which may be singular at the origin, while Oliva, Petitta and
Segura de Le\'on \cite{OPS2024} studied the regularizing role of
absorption terms. Filippucci and Zheng \cite{FZ2024} investigated
existence and nonexistence phenomena for weighted mean curvature
equations, including singular or degenerate weights.

The singular term in \eqref{P1} places the problem within the broad
theory of singular elliptic equations. Such equations require a
different treatment from standard superlinear problems because the
right-hand side is not defined at the homogeneous boundary value and
typically behaves like a negative power of the distance to
$\partial\Omega$; the corresponding boundary regularity and gradient
estimates go back to \cite{delPino,GuiLin} and, for quasilinear
operators, to \cite{GiacomoniSchindlerTakac}. A recent comprehensive account of this theory is
given by Oliva and Petitta \cite{OP2025}. In a related direction,
Balducci, Oliva and Petitta \cite{BOP2024} considered elliptic equations
where singular terms and nonlinear gradient contributions compete.
These developments emphasize that the simultaneous presence of
singularities and first-order terms is not merely a technical
perturbation: it affects the natural functional setting, compactness,
boundary estimates, and the construction of admissible barriers.

A second difficulty in \eqref{P1} is the convection term. Since
$f(x,u,\nabla u)$ depends explicitly on the gradient, the equation is in
general nonvariational. In particular, the usual energy associated with
the area functional does not directly encode the full right-hand side.
Sub--supersolution methods combined with truncation provide a flexible
alternative in this setting. This philosophy has recently proved useful
for quasilinear problems with convection and very large nonlinearities.
In \cite{AFM2024}, positive solutions were obtained for elliptic
equations with convection and variable exponents without an upper
restriction by introducing suitable contracted and truncated auxiliary
problems. Related ideas were developed for quasilinear systems in
\cite{AFMsys2024}. The present work adapts this mechanism to the
non-uniformly elliptic mean curvature operator and, at the same time,
incorporates the singular source $\lambda u^{-\alpha}$.

The interaction between these three features, the mean curvature
operator, the singular term, and the gradient-dependent
nonlinearity, is the main issue addressed here. The small-parameter
regime reveals the appropriate scale. Indeed, setting
\[
t_\lambda:=\lambda^{1/(1+\alpha)}
\]
gives the exact balance
\[
\lambda(t_\lambda w)^{-\alpha}
=t_\lambda w^{-\alpha}.
\]
Consequently, a supersolution of size $t_\lambda$ has its principal
elliptic contribution and its singular forcing at the same order.
Meanwhile, assumption \eqref{f} yields
\[
(t_\lambda w)^{r_1}=o(t_\lambda),
\qquad
|t_\lambda\nabla w|^{r_2}=o(t_\lambda)
\qquad(\lambda\to0^+),
\]
precisely because $r_1,r_2>1$. Thus these restrictions are not imposed
by a Sobolev critical threshold; rather, they express the perturbative
character of the zero- and first-order nonlinearities at the natural
singular scale. This observation is the key to the construction and
also explains why no Sobolev subcritical restriction on $r_1$ or
$r_2$ is needed in our argument.

Our approach can be summarized as follows. We first replace the mean
curvature coefficient outside the small-gradient region by a uniformly
elliptic coefficient which agrees with
$(1+|\nabla u|^2)^{-1/2}$ whenever $|\nabla u|<1$. The purely singular
auxiliary equation supplies a positive lower solution. For the upper
solution we use an explicit barrier at scale $t_\lambda$, which yields
an explicit smallness range for the parameter. We then contract the
gradient variable in $f$ and truncate the unknown between the ordered
lower and upper solutions. The resulting auxiliary problem can be
solved within a uniformly elliptic framework. Localization between the
barriers is recovered by comparison. Finally, after the rescaling
$u=t_\lambda w$, uniform $C^{1,\sigma}$ estimates show that the
gradient contraction is inactive for $\lambda$ small; the auxiliary
coefficient therefore coincides with the original mean curvature
coefficient. This closes the argument directly for \eqref{P1}.

The resulting conclusion is stronger than a qualitative existence
statement in the small-parameter regime: the proof identifies the
asymptotically correct size of the constructed solutions,
\[
\|u_\lambda\|_{C^1(\overline\Omega)}
=O\!\left(\lambda^{1/(1+\alpha)}\right).
\]
In this sense, the present result connects the modern theory of singular
mean curvature problems with the truncation/contraction techniques for
convection equations. The main contribution is the implementation of
this combined strategy for a singular prescribed mean curvature
equation with genuine gradient dependence, together with a quantitative
small-solution estimate and an explicit parameter regime produced by
the supersolution construction.

By a solution of \eqref{P1} we mean a function
$u\in C^1(\overline\Omega)$, positive in $\Omega$, such that
\[
\int_\Omega
\frac{\nabla u\cdot\nabla\varphi}
{\sqrt{1+|\nabla u|^2}}\,dx
=
\int_\Omega
\left(\frac{\lambda}{u^\alpha}
+f(x,u,\nabla u)\right)\varphi\,dx
\]
for every $\varphi\in H_0^1(\Omega)$. The restriction $0<\alpha<1$ is
particularly convenient here because the boundary behaviour of the
positive barriers makes the singular term compatible with this weak
formulation.

We can now state the main result.

\begin{theorem}\label{main result}
Assume $0<\alpha<1$ and that \eqref{f} holds with constant exponents
$r_1,r_2>1$. Then there exists $\lambda^*>0$ such that, for every
$\lambda\in(0,\lambda^*]$, problem \eqref{P1} possesses a positive weak
solution
\[
u_\lambda\in C^{1,\beta}(\overline\Omega)
\]
for some $\beta\in(0,1)$ independent of $\lambda$. Moreover, if
\[
t_\lambda:=\lambda^{1/(1+\alpha)},
\]
then there exists a constant $C>0$, independent of small $\lambda$, such
that
\[
\|u_\lambda\|_{C^1(\overline\Omega)}
\le Ct_\lambda.
\]
In particular,
\[
\|u_\lambda\|_{C^1(\overline\Omega)}
\longrightarrow0
\qquad\hbox{as }\lambda\to0^+.
\]
\end{theorem}

The paper is organized as follows. Section~\ref{auxiliary} introduces
the uniformly elliptic modification of the mean curvature operator and
the auxiliary singular problem. Section~\ref{S3} constructs ordered
sub- and supersolutions, with the supersolution chosen at the scale
$\lambda^{1/(1+\alpha)}$. Section~\ref{S4} introduces the order
truncation and the contraction in the gradient variable and establishes
existence for the truncated problem. Section~\ref{S5} derives the
uniform gradient estimate needed to remove the contraction. The final
section combines these ingredients to prove Theorem~\ref{main result}.

\section{Auxiliary results}\label{auxiliary}

In the course of the proof of Theorem \ref{main result}, we define an auxiliary quasilinear equation 
\begin{equation}\label{eq6}
	\left\{
	\begin{array}{lcl}
		-\operatorname{div}\left(a(|\nabla u|^2)\nabla\,u\right)  = \displaystyle\frac{\lambda}{u^{\alpha	}}  +f(x,u,\nabla u)\,\, &\mbox{ in }&\Omega, \\
		u>0 &\textup{ in }& \Omega\\
		u  = 0\,\, &\mbox{ on }& \partial \Omega.
	\end{array}
	\right.
\end{equation}
The function $a: [0,\infty)\rightarrow \mathbb{R}$ is defined by
\begin{equation}\label{eq7}
	a(v) = \left\{
	\begin{array}{rcl}
		\frac{1}{\sqrt{1+v}} & \, \, \mbox{ in } & [0,1), \\
		b(v-2)^2+c & \, \, \mbox{ in } & [1,2), \\
		c & \, \, \mbox{ in } & [2,+\infty)
	\end{array}
	\right.
\end{equation}
with $b=\frac{1}{8\sqrt{2}}$ and $c=\frac{7}{8\sqrt{2}}$. This function $a$ has been used in \cite{HabetsO}. Note that $a$ satisfies the structural ellipticity conditions of \cite{Lie} and gives rise to a uniformly elliptic operator. For bounded right-hand sides the global $C^{1,\sigma}$ theory of \cite{Lie} is therefore applicable. In the present problem, however, the term $u^{-\alpha}$ is unbounded at the boundary. The precise boundary-singular regularity statement needed below is recorded and justified in Section~\ref{S5}; see Lemma~\ref{lem:uniform-singular-C1beta}. In particular, notice that
\begin{equation}\label{ac}
	c \leq a(s) \leq 1, \forall s \in [0,+\infty). 
\end{equation}
Moreover
\begin{equation}\label{a0}
a(t) \in C^{1}(0,+\infty) \text { and } a(t)>0 \text { for } t>0.
\end{equation}

\begin{definition}\label{defsubsup}
Let us fix a value of $\lambda>0$. We say that a function $\overline\psi \in C^1(\overline{\Omega})$, positive in $\Omega$, is an upper solution of \eqref{eq6} if
$$
\overline\psi(x) \geq 0 \quad \text { on } \partial \Omega
$$
and
$$
\int_{\Omega} a\left(|\nabla \overline\psi|^2\right) \nabla \overline\psi\cdot \nabla w\, d x \geq \int_{\Omega} \Big(\frac{\lambda}{\overline\psi^{\alpha}}  +f(x,\overline\psi,\nabla \overline\psi)\Big)w\, d x
$$
for every $w \in H_0^1(\Omega)$, with $w(x) \geq 0$ a.e. in $\Omega$. Similarly, a lower solution $\underline\psi \in C^1(\overline{\Omega})$, positive in $\Omega$, is defined by reversing the two inequalities written above.
\end{definition}


Define the function 
\begin{equation}\label{Phi}
	\Phi(t)=\int_0^tsa(s^2)ds, \,\,\, t \in \mathbb{R}.
\end{equation}

%Aiming to prove Theorm \ref{main result} for $1<p<\frac{N+2}{N-2}$, we establish a precise bound estimate below. 
%
%\begin{proposition}\label{A_priori} Let $\gamma^*=\frac{2.2^{*}}{2^{*}-(p-1)}$, $C_*$ the embedding constant of $H_0^1(\Omega)$ in $L^{2^*}(\Omega)$. If $v_{\lambda}>0$ is a weak solution of the auxiliary problem (\ref{eq6}), satisfying	$\left\|v_{\lambda}\right\|_{H_{0}^{1}(\Omega)} \rightarrow 0$ as $\lambda \rightarrow 0$. Then, $v_{\lambda} \in L^{\infty}(\Omega)$ and
	%\begin{equation}\label{Linfity}
		%\|v_{\lambda}\|_{\infty} \leq C_5^{\sigma_1}\chi^{\sigma_2}\Theta(\lambda) 
	%\end{equation}
	%where $\chi=2^{*} / \gamma^{*}$,  $\sigma_1=\frac{1}{2}\left(\sum\limits_{i=1}^{\infty}\frac{1}{\chi^i}\right)$ and $\sigma_2=2\sum\limits_{i=1}^{\infty}\frac{i}{\chi^i}$ and
	%\[
	%\Theta(\lambda)=\max\{\lambda + \|v_{\lambda}\|_{H_{0}^{1}(\Omega)}^{p-1}, C_*\|v_{\lambda}\|_{H_{0}^{1}(\Omega)}\} 
	%\] 
	%satisfying $\Theta(\lambda)<1$ for $\lambda>0$ small enough. 
%\end{proposition}
%This result is of independent interest. Moreover, as a consequence, the solution obtained in Theorem \ref{main result} satisfy the $L^{\infty}$ estimate (\ref{Linfity}). 
%
%\begin{remark}
%Furthermore, we will prove that $v_{\lambda} \in C^{1,\sigma}(\overline{\Omega})$ with $0<\sigma <1$ and $v_{\lambda}$ is a solution of (\ref{P1})  satisfying
%\[
%\left\|v_{\lambda}\right\|_{C^{1}(\bar{\Omega})} \rightarrow 0 \mbox{ as } \lambda \rightarrow 0.
%\] 
%\end{remark}


The following result can be obtained by standard computations.
\begin{lemma}\label{cond.as} Let $a$ as defined in \eqref{eq7}, then
	 \begin{equation}\label{i}
	 	s a(s^2) \rightarrow 0 \mbox{ as } s \rightarrow 0, 
	 \end{equation}		\begin{equation}\label{ii}
	s a(s^2) \rightarrow \infty \mbox{ as } s \rightarrow \infty,
	\end{equation}
	\begin{equation}\label{iii}
	s \mapsto sa(s^2) \mbox{ is increasing in } (0, \infty),
	\end{equation} 
		\begin{equation}\label{iv}
		\frac{19}{40} \leq \frac{(s a(s^2))^{\prime}}{a(s^2)} \leq 1, \,\, s>0,
		\end{equation}
		\begin{equation}\label{radial}
		\nu_0:=\frac{19}{40\sqrt2}\le (s a(s^2))'=a(s^2)+2s^2a'(s^2)\le 1, \,\, s\ge0.
		\end{equation}
\end{lemma}
\begin{proof}
By the definition of $a$ in (\ref{eq7}), properties (\ref{i}) and (\ref{ii}) are easy.

Now, we will prove (\ref{iii}). By standard computations we can conclude that
\begin{equation}\label{eq7.3}
	a'(v) = \left\{
	\begin{array}{rcl}
		-\frac{1}{2(1+v)^{3/2}} & \, \, \mbox{ in } & [0,1), \\
		2b(v-2) & \, \, \mbox{ in } & [1,2), \\
		0 & \, \, \mbox{ in } & [2,+\infty).
	\end{array}
	\right.
\end{equation}
Hence, $a$ belongs to $C^1([0,+\infty))$. 

We will show that
\[
(sa(s^2))'=a(s^2)+2s^2a'(s^2)>0 \, \mbox{ for } \, s>0.
\]

Assume $s\in  (0,1)$. By (\ref{eq7}) and (\ref{eq7.3}) we obtain
\[
\begin{array}{rcl}
	a(s^2)+2s^2a'(s^2) &= & a(s^2)+2s^2\left(\frac{-1}{2(1+s^2)^{3/2}}\right) \\
	&= & \frac{1}{(1+s^2)^{1/2}}+2s^2\left(\frac{-1}{2(1+s^2)^{3/2}}\right)\\
	&=& \frac{1}{(1+s^2)^{1/2}}\left(1-\frac{s^2}{1+s^2}\right)=\frac{1}{(1+s^2)^{3/2}}\geq \frac{1}{2\sqrt2}>\nu_0.
\end{array}
\]

For $s\in  [1,\sqrt{2})$ put $v=s^2\in[1,2)$. By (\ref{eq7}), (\ref{eq7.3}) and $c=7b$ we obtain
\begin{equation}\label{middle}
\begin{array}{rcl}
	a(v)+2va'(v) &= & b(v-2)^2+c+4bv(v-2) \\
	&= & b\left(5v^2-12v+11\right)\\
	&=& 5b\left(v-\frac65\right)^2+\frac{19}{5}\,b\\
	&\geq& \frac{19}{5}\,b=\frac{19}{40\sqrt2}=\nu_0>0.
\end{array}
\end{equation}

Assume $s\in  [\sqrt{2},\infty)$. Using (\ref{eq7}) and (\ref{eq7.3}) we obtain
\[
a(s^2)+2s^2a'(s^2) = c+0=c=\frac{7}{8\sqrt2}>\nu_0.
\]
Therefore, $s\,a(s^2)$ is increasing in $(0, \infty)$. Since $a'\le0$ and $a\le1$, we also have $a(s^2)+2s^2a'(s^2)\le a(s^2)\le1$, which proves \eqref{radial}.

Next, we prove (\ref{iv}).

By standard computations
\[
\frac{(sa(s^2))'}{a(s^2)}=\frac{a(s^2)+2s^2a'(s^2)}{a(s^2)}=1+\frac{2s^2a'(s^2)}{a(s^2)}.
\]
Assume $s\in  (0,1)$. By (\ref{eq7}) and (\ref{eq7.3}) we obtain
\[
\begin{array}{rcl}
	\displaystyle\frac{2s^2a'(s^2)}{a(s^2)} &= & \displaystyle\frac{2s^2\left(-\frac{1}{2(1+s^2)^{3/2}}\right)}{\frac{1}{(1+s^2)^{1/2}}} \\
	&= &\displaystyle \frac{-2s^2}{2(1+s^2)^{3/2}}(1+s^2)^{1/2}\\
	&=&\displaystyle \frac{-s^2}{1+s^2}.
\end{array}
\]
Hence,
\[
1+ \displaystyle\frac{2s^2a'(s^2)}{a(s^2)} = 1 - \frac{s^2}{1+s^2}=\frac{1}{1+s^2}>\frac{1}{2}
\]
and moreover
\[
1+ \displaystyle\frac{2s^2a'(s^2)}{a(s^2)} \leq 1.
\]
Therefore,
\begin{equation}\label{ineq1}
	\frac{1}{2}<\frac{(sa(s^2))'}{a(s^2)}\leq 1.
\end{equation}

Assume $s\in  [1,\sqrt{2})$, in particular $s^2-2<0$. Using (\ref{eq7}) and (\ref{eq7.3}) we obtain
\[
\begin{array}{rcl}
	\displaystyle\frac{2s^2a'(s^2)}{a(s^2)} &= & \displaystyle\frac{4bs^2(s^2-2)}{b(s^2-2)^2+c}.
\end{array}
\]
Hence,
\[
1+\displaystyle\frac{2s^2a'(s^2)}{a(s^2)}=1+ \frac{4bs^2(s^2-2)}{b(s^2-2)^2+c}<1.
\]
Moreover, since $b(s^2-2)^2+c\le b+c=8b$ for $s^2\in[1,2)$, \eqref{middle} gives
\[
	1+\displaystyle\frac{2s^2a'(s^2)}{a(s^2)} =  \displaystyle\frac{a(s^2)+ 2s^2a'(s^2)}{b(s^2-2)^2+c}
	\geq  \displaystyle\frac{\frac{19}{5}\,b}{8b}=\frac{19}{40}.
\]
Then
\begin{equation}\label{ineq2}
	\frac{19}{40}\le\frac{(sa(s^2))'}{a(s^2)}< 1.
\end{equation}

Assume $s\in  [\sqrt{2},+\infty)$. By (\ref{eq7}) and (\ref{eq7.3}) we obtain
\begin{equation}\label{ineq3}
	\frac{(sa(s^2))'}{a(s^2)}=1+\displaystyle\frac{2s^2a'(s^2)}{a(s^2)}=1.
\end{equation}

As a consequence of (\ref{ineq1}), (\ref{ineq2}) and (\ref{ineq3}), we get
$$
\frac{19}{40} \leq \frac{(s a(s^2))^{\prime}}{a(s^2)} \leq 1, \,\, s>0.
$$
%with
%\[
%\ell=\frac{3}{2} \text{ and } m=2.
%\]
\end{proof}

\begin{remark}\label{rem:sharp}
The constant $\frac12$ cannot replace $\frac{19}{40}$ in \eqref{iv}. Indeed,
for $v=s^2\in[1,2)$ one has
$\frac{(sa(s^2))'}{a(s^2)}=\frac{5v^2-12v+11}{v^2-4v+11}$, whose minimum
on $[1,2]$ is attained at $v=(11-\sqrt{77})/2\approx1.113$ and is
approximately $0.4929<\frac12$. Only the positivity of the lower bound
is used in the sequel.
\end{remark}

\begin{remark}\label{rem:DA}
For $\mathcal A(\xi):=a(|\xi|^2)\xi$ one has
$D\mathcal A(\xi)=a(|\xi|^2)I+2a'(|\xi|^2)\xi\otimes\xi$. Its eigenvalues
are $a(|\xi|^2)\in[c,1]$ (in the directions orthogonal to $\xi$) and
$a(|\xi|^2)+2|\xi|^2a'(|\xi|^2)\in[\nu_0,1]$ (in the direction of $\xi$),
by \eqref{ac} and \eqref{radial}. Since $\nu_0<c$, it follows that
\begin{equation}\label{DA-bounds}
\nu_0|\eta|^2\le D\mathcal A(\xi)\eta\cdot\eta\le|\eta|^2,
\qquad \xi,\eta\in\mathbb R^N,
\end{equation}
and consequently
\begin{equation}\label{strongmono}
\bigl(\mathcal A(\xi)-\mathcal A(\zeta)\bigr)\cdot(\xi-\zeta)\ge\nu_0|\xi-\zeta|^2,
\qquad \xi,\zeta\in\mathbb R^N.
\end{equation}
The same bounds hold, with the same constants, for
$\mathcal A_t(\xi):=a(t^2|\xi|^2)\xi$ for every $t>0$, since
$D\mathcal A_t(\xi)=D\mathcal A(t\xi)$.
\end{remark}

An immediate conclusion is the following one.

\begin{lemma}\label{cor.as}
	Let $\phi(s)=a(s^2)$, properties (\ref{i})-(\ref{iv}) of Lemma \ref{cond.as} imply
\begin{equation}\label{phi1}
s \phi(s) \rightarrow 0 \mbox{ as } s \rightarrow 0, 
\end{equation}
\begin{equation}\label{phi2}
	s \phi(s) \rightarrow \infty \mbox{ as } s \rightarrow \infty,
\end{equation}
 \begin{equation}\label{phi3}
 s \mapsto s \phi(s) \mbox{ is strictly increasing in } (0, \infty), 
 \end{equation}
\begin{equation}\label{phi4}
			\frac{19}{40}\leq \frac{(s \phi(s))^{\prime}}{\phi(s)} \leq 1, \,\, s>0.
\end{equation}		
\end{lemma}


For the reader's convenience we recall some facts on Orlicz--Sobolev spaces; see \cite{RR}.
	
\noindent By (\ref{a0}) and Lemma \ref{cor.as}-(\ref{phi4}), the Orlicz space $L_{\Phi}(\Omega)$ coincides with the set (of equivalence classes) of measurable functions $u: \Omega \rightarrow \mathbb{R}$ such that
	\begin{equation}\label{Orl}
		\int_{\Omega} \Phi(|u|) dx<+\infty.
	\end{equation}
The usual norm on $L_{\Phi}(\Omega)$ is
	$$
	\|u\|_{\Phi}=\inf \left\{\mu>0: \int_{\Omega} \Phi\left(\frac{u(x)}{\mu}\right) d x \leq 1\right\}
	$$
	called the Luxemburg norm, and the norm on $W^{1, \Phi}(\Omega)$ is
	$$
	\|u\|_{1, \Phi}=\|u\|_{\Phi}+\sum_{i=1}^{N}\left\|\frac{\partial u}{\partial x_{i}}\right\|_{\Phi},
	$$
	known as Orlicz-Sobolev norm, while $W_{0}^{1, \Phi}(\Omega)$ stands for the closure of $C_{0}^{\infty}(\Omega)$ with respect to the norm in $W^{1, \Phi}(\Omega)$. %Also, by (?????) the Banach spaces $L_{\Phi}(\Omega)$ and $W^{1, \Phi}(\Omega)$ are separable and reflexive. 
	
	
\noindent	By (\ref{ac}) and (\ref{Phi}) we obtain 
	\begin{equation}\label{eq26}
		\frac{c}{2}t^2\leq \Phi(t)\leq \frac{1}{2}t^2, \,\, \mbox{ for } \,\, t\geq 0.
	\end{equation}
	Hence,
	\begin{equation}\label{eq32}
		\frac{c}{2\mu^2}\|u\|_{H_0^1(\Omega)}^2\leq \int_{\Omega}\Phi\left(\frac{|\nabla u|}{\mu}\right)dx \leq \frac{1}{2\mu^2}\|u\|_{H_0^1(\Omega)}^2, \,\,\, \forall u \in  W^{1,2}_0(\Omega),\ \mu>0.
	\end{equation}
	
\noindent	By (\ref{eq26}) and the fact that $c<1$, we obtain
	\begin{equation}\label{eq30}
		\frac{1}{2}t^2\leq \frac{1}{c}\Phi(t) \leq \Phi\left(\frac{1}{c}t\right),
	\end{equation}
	see \cite[Lemma 3.2.2]{KOF}. 
	
\noindent	By (\ref{eq26}) and (\ref{eq32})
	\[
	L_{\Phi}(\Omega)=L^2(\Omega)
	\]
and
	\begin{equation}\label{eq31}
		W_{0}^{1, \Phi}(\Omega)=H^1_0(\Omega).
	\end{equation}


Now, $t^2\phi(t)=t\Phi'(t)$ and, by \eqref{phi4},
$\frac{19}{40}\Phi'(t)\le t\Phi''(t)=t\,(t\phi(t))'\le\Phi'(t)$ for $t>0$.
Hence $\frac{59}{40}\Phi'\le (t\Phi')'\le2\Phi'$, and integrating on
$(0,t)$ (note that $t\Phi'(t)\to0$ as $t\to0$) we obtain
\begin{equation}\label{phi5}
		\frac{59}{40}\leq \frac{t^2\phi(t)}{\Phi(t)} \leq 2, \,\, t>0.
\end{equation}
Defining
\begin{equation}\label{infphi}
	\underline{\phi}:=\inf _{t>0} \frac{t^2\phi(t)}{\Phi(t)} \mbox{ and } \overline{\phi}:=\sup _{t>0} \frac{t^2 \phi(t)}{\Phi(t)},
\end{equation}
we obtain that
\begin{equation}\label{underoverphi}
	\frac{59}{40}\leq \underline{\phi} \,\,\mbox{ and } \,\, \overline{\phi}\leq 2.
\end{equation}


\begin{remark}\label{rema2}
	By (\ref{phi5}), the function $\Phi(t)$ satisfies the $\Delta_{2}$-condition, that is
	$$
	\exists k>0,\,\,	\Phi(2 t) \leq k \Phi(t), \quad t>0.
	$$
\end{remark}


Firstly, we consider the following auxiliary problem
\begin{equation}\label{eq9}
	\left\{
	\begin{array}{rcl}
		-\operatorname{div}\left(a(|\nabla u|^2)\nabla\,u\right)   = \displaystyle\frac{\lambda}{u^{\alpha	}}&\mbox{ in }& \Omega, \\
		u>0 &\textup{ in }& \Omega\\
		u = 0&\mbox{ on }& \partial \Omega.
	\end{array}
	\right.
\end{equation} 
\begin{lemma}\label{le2}
	For any constant $\lambda>0$ the problem (\ref{eq9})
	admits a solution $u_{0,\lambda}\in C^1(\overline{\Omega})$ and $\partial u_{0,\lambda}/\partial \eta <0$ on $\partial\Omega$.
\end{lemma}
\begin{proof}
For $\varepsilon>0$ consider the regularized problem
\begin{equation}\label{eq9eps}
\left\{
\begin{array}{rcl}
-\operatorname{div}\!\left(a(|\nabla u|^2)\nabla u\right)
&=&\displaystyle\frac{\lambda}{(u+\varepsilon)^\alpha}
\quad\hbox{in }\Omega,\\
u&=&0\quad\hbox{on }\partial\Omega .
\end{array}\right.
\end{equation}
Since $c\le a\le1$ and $s\mapsto s a(s^2)$ is strictly increasing, the
operator
\[
A(u)=-\operatorname{div}\!\left(a(|\nabla u|^2)\nabla u\right)
\]
is bounded, coercive, strictly monotone and of type $(S_+)$ from
$H_0^1(\Omega)$ into $H^{-1}(\Omega)$. The right-hand side of
\eqref{eq9eps} is bounded. Hence the Browder--Minty theorem gives a unique
$u_\varepsilon\in H_0^1(\Omega)$. Testing with $u_\varepsilon^-$ gives
$u_\varepsilon\ge0$; by \cite{Lie}, $u_\varepsilon\in C^{1,\sigma}(\overline\Omega)$,
and the strong maximum principle \cite{Mont.} yields
$u_\varepsilon>0$ in $\Omega$.

If $0<\varepsilon_1<\varepsilon_2$, comparison gives
$u_{\varepsilon_1}\ge u_{\varepsilon_2}$. Fixing one regularized solution
as a positive lower barrier and using the standard boundary barrier
$d(x)=\operatorname{dist}(x,\partial\Omega)$, one obtains constants
$m_1,m_2>0$, independent of $\varepsilon\in(0,1]$, such that
\[
 m_1d(x)\le u_\varepsilon(x)\le m_2
 \qquad\hbox{in }\Omega.
\]
Consequently
\[
0\le \frac{\lambda}{(u_\varepsilon+\varepsilon)^\alpha}
\le C\,d(x)^{-\alpha}.
\]
Because $0<\alpha<1$, Hardy's inequality shows that
$d^{-\alpha}$ defines an element of $H^{-1}(\Omega)$:
\[
\int_\Omega d^{-\alpha}|\varphi|
\le
\left(\int_\Omega d^{2(1-\alpha)}\right)^{1/2}
\left\|\frac{\varphi}{d}\right\|_2
\le C\|\nabla\varphi\|_2 .
\]
The energy estimate and monotonicity therefore give, along a subsequence,
$u_\varepsilon\to u_{0,\lambda}$ strongly in $H_0^1(\Omega)$ and a.e. in
$\Omega$. Passing to the limit in \eqref{eq9eps} yields
\[
-\operatorname{div}\!\left(a(|\nabla u_{0,\lambda}|^2)
\nabla u_{0,\lambda}\right)=\lambda u_{0,\lambda}^{-\alpha}.
\]
The boundary estimate above persists in the limit. Since the right-hand
side is bounded by $Cd^{-\alpha}$ with $0<\alpha<1$, the boundary regularity
for singular quasilinear equations \cite[Theorem B.1]{GiacomoniSchindlerTakac}
(see also \cite{GuiLin,delPino} for the semilinear case) gives
$u_{0,\lambda}\in C^{1,\sigma}(\overline\Omega)$ for some
$\sigma\in(0,1)$; the strong maximum principle and Hopf boundary lemma \cite{Mont.} give
$u_{0,\lambda}>0$ in $\Omega$ and
$\partial u_{0,\lambda}/\partial\eta<0$ on $\partial\Omega$.
\end{proof}

\begin{definition}
	A function $u \in W^{1, \Phi}(\Omega)$, positive in $\Omega$, is a subsolution of (\ref{eq9}) if
	$$
	\int_{\Omega} \phi(|\nabla u|) \nabla u \nabla \varphi d x \leq \int_{\Omega} \displaystyle\frac{\lambda}{u^{\alpha	}} \varphi d x, \,\, \forall \varphi \in W_{0}^{1, \Phi}(\Omega), \,\, \varphi \geq 0.
	$$
	A function $u \in W^{1, \Phi}(\Omega)$, positive in $\Omega$, is a supersolution of (\ref{eq9}) if
	$$
	\int_{\Omega} \phi(|\nabla u|) \nabla u \nabla \varphi d x \geq \int_{\Omega} \displaystyle\frac{\lambda}{u^{\alpha	}}\varphi d x, \,\,  \forall  \varphi \in W_{0}^{1, \Phi}(\Omega), \,\, \varphi \geq 0.
	$$
	A function $u \in W^{1, \Phi}(\Omega)$, positive in $\Omega$, is a solution of (\ref{eq9}) if
	$$
	\int_{\Omega} \phi(|\nabla u|) \nabla u \nabla \varphi d x = \int_{\Omega} \displaystyle\frac{\lambda}{u^{\alpha	}}\varphi d x,  \,\,   \forall   \varphi \in W_{0}^{1,\Phi}(\Omega).
	$$
\end{definition}
The following comparison principle is a simple version, adapted to the
present setting, of \cite[Theorem 1.2]{GCS}.
\begin{lemma} [Comparison Principle] \label{comp.lemma}
	Let $u_{1}, u_{2} \in W^{1, \Phi}(\Omega)$ be a subsolution and a supersolution of (\ref{eq9}), respectively, with $u_1,u_2>0$ in $\Omega$, $(u_1-u_2)^+\in W_0^{1,\Phi}(\Omega)$, and such that $u_i^{-\alpha}(u_1-u_2)^+\in L^1(\Omega)$, $i=1,2$. Then $u_{1} \leq u_{2}$ a.e. in $\Omega$.
\end{lemma}
\begin{proof}
Testing the two inequalities with $\varphi=(u_1-u_2)^+\ge0$ and subtracting, we obtain
\[
\int_{\{u_1>u_2\}}\bigl(\mathcal A(\nabla u_1)-\mathcal A(\nabla u_2)\bigr)\cdot\nabla(u_1-u_2)\,dx
\le\lambda\int_{\{u_1>u_2\}}\bigl(u_1^{-\alpha}-u_2^{-\alpha}\bigr)(u_1-u_2)\,dx\le0,
\]
because $s\mapsto s^{-\alpha}$ is decreasing. By \eqref{strongmono},
$\nu_0\|\nabla(u_1-u_2)^+\|_2^2\le0$, hence $(u_1-u_2)^+=0$.
\end{proof}

\section{Sub--supersolutions for problem \eqref{eq6}}\label{S3}

We construct the supersolution at the scale dictated by the singular
term. Put
\[
 t_\lambda:=\lambda^{\frac1{1+\alpha}},\qquad
 \lambda=t_\lambda^{1+\alpha}.
\]
Choose $x_0\in\mathbb R^N\setminus\overline\Omega$ and set
\[
 R_\Omega:=\max_{x\in\overline\Omega}|x-x_0|,\qquad
 r_\Omega:=\min_{x\in\overline\Omega}|x-x_0|>0.
\]
Fix $M_0:=R_\Omega^2+1$ and $\vartheta(x):=M_0-|x-x_0|^2$. Then
\begin{equation}\label{beta-bounds}
1\le\vartheta(x)\le B_0:=M_0-r_\Omega^2\le M_0,\qquad
|\nabla\vartheta(x)|=2|x-x_0|\le G:=2R_\Omega.
\end{equation}
Define
\begin{equation}\label{cstar}
 C_*:=\frac{N}{\sqrt2},\qquad \delta_*:=C_*-1>0,
\end{equation}
where $\delta_*>0$ because $N\ge2$, and
\begin{equation}\label{Lambda-explicit}
\boxed{\Lambda:=\min\left\{
1,G^{-(1+\alpha)},
\left(\frac{\delta_*}{2k_1B_0^{r_1}}\right)^{\frac{1+\alpha}{r_1-1}},
\left(\frac{\delta_*}{2k_2G^{r_2}}\right)^{\frac{1+\alpha}{r_2-1}}
\right\}.}
\end{equation}
Since $r_1,r_2>1$, one has $\Lambda>0$.

For $0<\lambda\le\Lambda$ set
\begin{equation}\label{subsuper-def}
 \underline u_\lambda:=u_{0,\lambda},\qquad
 \overline u_\lambda:=t_\lambda\vartheta.
\end{equation}

\begin{theorem}\label{Th2}
Assume $0<\alpha<1$ and \eqref{f} with constant exponents $r_1,r_2>1$.
For every $\lambda\in(0,\Lambda]$, $\underline u_\lambda$ is a subsolution
of \eqref{eq6}, $\overline u_\lambda$ is a supersolution of \eqref{eq6}, and
\[
0<\underline u_\lambda\le\overline u_\lambda\quad\hbox{in }\Omega.
\]
Moreover,
\begin{equation}\label{super-C1}
\|\overline u_\lambda\|_{C^1(\overline\Omega)}
\le(B_0+G)\lambda^{\frac1{1+\alpha}},
\qquad
\|\nabla\overline u_\lambda\|_\infty\le Gt_\lambda\le1.
\end{equation}
\end{theorem}

\begin{proof}
Since $u_{0,\lambda}$ solves \eqref{eq9} and $f\ge0$, it is immediately a
subsolution of \eqref{eq6}.

Set $t=t_\lambda$. By \eqref{Lambda-explicit}, $t\le G^{-1}$, whence
$|\nabla(t\vartheta)|\le Gt\le1$. Thus the coefficient $a$ coincides along
$t\vartheta$ with the mean-curvature coefficient. Writing $y=x-x_0$, direct
calculation gives
\begin{align}
-&\operatorname{div}\left(\frac{\nabla(t\vartheta)}
{\sqrt{1+|\nabla(t\vartheta)|^2}}\right)\notag\\
&=\frac{2Nt}{\sqrt{1+4t^2|y|^2}}
-\frac{8t^3|y|^2}{(1+4t^2|y|^2)^{3/2}}\notag\\
&=\frac{2t\,[N+4(N-1)t^2|y|^2]}
{(1+4t^2|y|^2)^{3/2}}.\label{exact-super-op}
\end{align}
Consequently,
\begin{equation}\label{lhs-super-lower}
-\operatorname{div}\left(a(|\nabla\overline u_\lambda|^2)
\nabla\overline u_\lambda\right)
\ge\frac{2Nt}{(1+4t^2R_\Omega^2)^{3/2}}
\ge\frac{N}{\sqrt2}t=C_*t.
\end{equation}

Since $\lambda=t^{1+\alpha}$ and $\vartheta\ge1$,
\begin{equation}\label{sing-super}
\frac{\lambda}{\overline u_\lambda^\alpha}
=t\vartheta^{-\alpha}\le t.
\end{equation}
Also, by \eqref{f} and \eqref{beta-bounds},
\begin{equation}\label{f-super}
f(x,\overline u_\lambda,\nabla\overline u_\lambda)
\le k_1B_0^{r_1}t^{r_1}+k_2G^{r_2}t^{r_2}.
\end{equation}
The definition of $\Lambda$ yields
\[
k_1B_0^{r_1}t^{r_1}\le\frac{\delta_*}{2}t,
\qquad
k_2G^{r_2}t^{r_2}\le\frac{\delta_*}{2}t.
\]
Hence
\[
\frac{\lambda}{\overline u_\lambda^\alpha}
+f(x,\overline u_\lambda,\nabla\overline u_\lambda)
\le(1+\delta_*)t=C_*t.
\]
Together with \eqref{lhs-super-lower}, this proves that
$\overline u_\lambda$ is a supersolution of \eqref{eq6}.

Since it is also a supersolution of the singular problem \eqref{eq9}, while
$u_{0,\lambda}$ is a solution of \eqref{eq9}, Lemma~\ref{comp.lemma}
(applied with $u_1=u_{0,\lambda}$, $u_2=\overline u_\lambda$; note that
$(u_{0,\lambda}-\overline u_\lambda)^+$ has compact support in $\Omega$,
since $\overline u_\lambda\ge t_\lambda>0$ on $\overline\Omega$) gives
\[
u_{0,\lambda}\le\overline u_\lambda\quad\hbox{in }\Omega,
\]
because $u_{0,\lambda}=0\le\overline u_\lambda$ on $\partial\Omega$.
Thus the pair is ordered. Finally \eqref{beta-bounds} gives
\eqref{super-C1}.
\end{proof}


\section{Gradient contraction and truncated problem}\label{S4}

The argument below follows the contraction--truncation strategy used for
convection terms in \cite{AFM2024}.  The essential point is that the
gradient dependence is first contracted to a bounded ball and only then
the order truncation associated with the sub--supersolution pair is
introduced.  This separates existence from the a priori estimate needed
to remove the contraction.

Put
\[
t_\lambda:=\lambda^{1/(1+\alpha)}.
\]
For $R>0$ define the radial retraction $\pi_R:\mathbb R^N\to
\overline B_R(0)$ by
\[
\pi_R(\xi)=
\begin{cases}
\xi,&|\xi|\le R,\\[1mm]
R\,\dfrac{\xi}{|\xi|},&|\xi|>R,
\end{cases}
\]
and set
\[
f_R(x,s,\xi):=f(x,s,\pi_R(\xi)).
\]
Let
\[
T_\lambda(s)(x):=
\begin{cases}
\underline u_\lambda(x),&s<\underline u_\lambda(x),\\
s,&\underline u_\lambda(x)\le s\le\overline u_\lambda(x),\\
\overline u_\lambda(x),&s>\overline u_\lambda(x).
\end{cases}
\]
For $u\in H_0^1(\Omega)$ we write simply $T_\lambda u(x)=T_\lambda(u(x))(x)$.

We consider
\begin{equation}\label{PRT}
\left\{
\begin{array}{rcl}
-\operatorname{div}\left(a(|\nabla u|^2)\nabla u\right)
&=&\displaystyle
\frac{\lambda}{(T_\lambda u)^\alpha}
+f_R\!\left(x,T_\lambda u,\nabla(T_\lambda u)\right)
\quad\hbox{in }\Omega,\\
u&=&0\quad\hbox{on }\partial\Omega .
\end{array}
\right.
\end{equation}

\begin{lemma}\label{singularfunctional}
There exists $C_\lambda>0$ such that
\[
\left|\int_\Omega (T_\lambda u)^{-\alpha}\varphi\,dx\right|
\le C_\lambda\|\varphi\|_{H_0^1(\Omega)}
\]
for all $u,\varphi\in H_0^1(\Omega)$.
\end{lemma}

\begin{proof}
Since $T_\lambda u\ge\underline u_\lambda$ and the Hopf estimate for
$\underline u_\lambda$ gives
$\underline u_\lambda(x)\ge c_\lambda d(x)$, we have
$(T_\lambda u)^{-\alpha}\le c_\lambda^{-\alpha}d(x)^{-\alpha}$.
Therefore
\[
\int_\Omega (T_\lambda u)^{-\alpha}|\varphi|
\le C_\lambda\int_\Omega d^{1-\alpha}\frac{|\varphi|}{d}
\le C_\lambda\|d^{1-\alpha}\|_2
\left\|\frac{\varphi}{d}\right\|_2.
\]
Hardy's inequality completes the proof.
\end{proof}

\begin{proposition}\label{existtruncated}
For every sufficiently small $\lambda>0$ and every $R>0$, problem
\eqref{PRT} possesses a weak solution $u_{\lambda,R}\in H_0^1(\Omega)$.
\end{proposition}

\begin{proof}
Define $\mathcal A_{\lambda,R}:H_0^1(\Omega)\to H^{-1}(\Omega)$ by
\[
\begin{aligned}
\langle\mathcal A_{\lambda,R}(u),v\rangle
={}&\int_\Omega a(|\nabla u|^2)\nabla u\cdot\nabla v\,dx\\
&-\int_\Omega\left[
\frac{\lambda}{(T_\lambda u)^\alpha}
+f_R(x,T_\lambda u,\nabla(T_\lambda u))
\right]v\,dx.
\end{aligned}
\]
The principal part is bounded, coercive, strictly monotone and of type
$(S_+)$.  By Lemma \ref{singularfunctional}, the singular term is bounded
in $H^{-1}(\Omega)$.  Moreover, since
$0\le T_\lambda u\le\|\overline u_\lambda\|_\infty$ and
$|\pi_R(\nabla T_\lambda u)|\le R$, assumption \eqref{f} gives
\[
0\le f_R(x,T_\lambda u,\nabla T_\lambda u)
\le k_1\|\overline u_\lambda\|_\infty^{r_1}+k_2R^{r_2}.
\]
Thus the lower-order part maps bounded sets into bounded subsets of
$H^{-1}(\Omega)$.

Let $u_n\rightharpoonup u$ in $H_0^1(\Omega)$ with
$\limsup\langle\mathcal A_{\lambda,R}(u_n),u_n-u\rangle\le0$. By the
compact embedding $H_0^1(\Omega)\hookrightarrow L^2(\Omega)$, $u_n\to u$ in
$L^2(\Omega)$. Since the contracted convection term is bounded in
$L^\infty(\Omega)$ uniformly in $n$,
\[
\left|\int_\Omega f_R(x,T_\lambda u_n,\nabla(T_\lambda u_n))(u_n-u)\,dx\right|
\le C\|u_n-u\|_1\to0 .
\]
For the singular term, let $\delta>0$ and split $\Omega$ into
$\{d<\delta\}$ and $\{d\ge\delta\}$. Using
$(T_\lambda u_n)^{-\alpha}\le c_\lambda^{-\alpha}d^{-\alpha}$ and Hardy's
inequality as in Lemma~\ref{singularfunctional},
\[
\int_\Omega (T_\lambda u_n)^{-\alpha}|u_n-u|\,dx
\le C\delta^{1-\alpha}\|\nabla(u_n-u)\|_2+C\delta^{-\alpha}\|u_n-u\|_1 .
\]
Since $(u_n)$ is bounded in $H_0^1(\Omega)$, letting first $n\to\infty$
and then $\delta\to0$ shows that this integral tends to zero. Therefore
\[
\limsup_{n\to\infty}
\int_\Omega a(|\nabla u_n|^2)\nabla u_n\cdot
(\nabla u_n-\nabla u)\,dx\le0.
\]
The $(S_+)$ property yields $u_n\to u$ strongly in $H_0^1(\Omega)$.
Hence $\mathcal A_{\lambda,R}$ is pseudomonotone.

Finally, using $a\ge c>0$, Poincar\'e's inequality and the boundedness
established above,
\[
\langle\mathcal A_{\lambda,R}(u),u\rangle
\ge c\|\nabla u\|_2^2-C_{\lambda,R}\|\nabla u\|_2,
\]
so $\mathcal A_{\lambda,R}$ is coercive.  Surjectivity of bounded
coercive pseudomonotone operators gives
$\mathcal A_{\lambda,R}(u_{\lambda,R})=0$.
\end{proof}

\begin{proposition}\label{location}
Let $0<\lambda\le\Lambda$ and $R\ge Gt_\lambda$. Then every solution
supplied by Proposition \ref{existtruncated} satisfies
\[
\underline u_\lambda\le u_{\lambda,R}\le\overline u_\lambda
\qquad\hbox{a.e. in }\Omega.
\]
Consequently $T_\lambda u_{\lambda,R}=u_{\lambda,R}$ and
$u_{\lambda,R}$ solves
\begin{equation}\label{PR}
-\operatorname{div}\left(a(|\nabla u|^2)\nabla u\right)
=\lambda u^{-\alpha}+f_R(x,u,\nabla u).
\end{equation}
\end{proposition}

\begin{proof}
Set $u=u_{\lambda,R}$.  Test the equation with
$(u-\overline u_\lambda)^+$.  On the set
$\{u>\overline u_\lambda\}$ one has
$T_\lambda u=\overline u_\lambda$ and
$\nabla(T_\lambda u)=\nabla\overline u_\lambda$ a.e.  Since
$|\nabla\overline u_\lambda|\le Gt_\lambda\le R$ by \eqref{super-C1}, we have
$\pi_R(\nabla\overline u_\lambda)=\nabla\overline u_\lambda$, so that
$f_R(x,T_\lambda u,\nabla(T_\lambda u))=f(x,\overline u_\lambda,\nabla\overline u_\lambda)$
on this set.  (This is the only place where the restriction
$R\ge Gt_\lambda$ is used: no monotonicity of $f$ in $\xi$ is assumed.)
Subtracting the supersolution inequality for $\overline u_\lambda$ gives
\[
\int_{\{u>\overline u_\lambda\}}
\left[
a(|\nabla u|^2)\nabla u
-a(|\nabla\overline u_\lambda|^2)\nabla\overline u_\lambda
\right]\cdot
(\nabla u-\nabla\overline u_\lambda)\,dx\le0.
\]
By \eqref{strongmono}, $(u-\overline u_\lambda)^+=0$.  Testing
similarly with $(\underline u_\lambda-u)^+$, using the subsolution
inequality and $f_R\ge0$, gives $\underline u_\lambda\le u$.  Hence
$T_\lambda u=u$, and \eqref{PR} follows.
\end{proof}

\section{Removal of the gradient contraction}\label{S5}

We record the uniform estimate needed to pass from \eqref{PR} to the
original equation.  It is the singular analogue, for the present
uniformly elliptic mean-curvature perturbation, of the a priori gradient
estimate used in the contraction argument of \cite{AFM2024}.


We next record the two boundary estimates needed in the rescaled
argument.  They isolate the only delicate regularity point caused by
the singular term.  The second lemma uses precisely the singular
boundary regularity result in \cite[Theorem B.1]{GiacomoniSchindlerTakac},
which is an adaptation of Lieberman's boundary regularity method
\cite{Lie} to right-hand sides with a boundary singularity.

\begin{lemma}[Uniform linear boundary estimate]
\label{lem:uniform-linear-boundary}
Let $\Omega\subset\mathbb R^N$ be a bounded $C^2$ domain and
$0<\alpha<1$.  Let $t_\lambda\to0$ and define
\[
 \mathcal A_\lambda(\xi):=a(t_\lambda^2|\xi|^2)\xi .
\]
Assume that there exist constants $\nu,L>0$, independent of
$\lambda$, such that
\begin{equation}\label{eq:A-uniform-ellipticity}
 \nu|\eta|^2
 \le D\mathcal A_\lambda(\xi)\eta\cdot\eta,
 \qquad
 |D\mathcal A_\lambda(\xi)|\le L
\end{equation}
for every $\xi,\eta\in\mathbb R^N$.  Suppose that
$w_\lambda\in H_0^1(\Omega)\cap L^\infty(\Omega)$ is positive and
satisfies
\begin{equation}\label{eq:abstract-scaled}
 -\operatorname{div}\mathcal A_\lambda(\nabla w_\lambda)
 =
 w_\lambda^{-\alpha}+h_\lambda(x)
 \quad\hbox{in }\Omega,
 \qquad
 w_\lambda=0\quad\hbox{on }\partial\Omega,
\end{equation}
where, uniformly in $\lambda$,
\begin{equation}\label{eq:uniform-preboundary}
 0\le h_\lambda\le H,\qquad
 \|w_\lambda\|_\infty\le M,\qquad
 w_\lambda(x)\ge c_0d(x).
\end{equation}
Then there exists $C_d>0$, independent of $\lambda$, such that
\begin{equation}\label{eq:two-sided-distance}
 c_0d(x)\le w_\lambda(x)\le C_dd(x)
 \qquad (x\in\Omega).
\end{equation}
\end{lemma}

\begin{proof}
Only the upper estimate has to be proved.  Since $\partial\Omega$ is
$C^2$, there exists $\rho_0>0$ such that
\[
 \Omega_{\rho_0}:=\{x\in\Omega:0<d(x)<\rho_0\}
\]
is a tubular neighborhood of the boundary,
$d\in C^2(\overline{\Omega_{\rho_0}})$, $|\nabla d|=1$, and
$|D^2d|\le K_d$ there.

Fix positive constants $A,E$ and put
\[
 q(s):=As-Es^{2-\alpha},\qquad
 \Psi(x):=q(d(x)).
\]
Then
\[
 q'(s)=A-E(2-\alpha)s^{1-\alpha},\qquad
 q''(s)=-E(2-\alpha)(1-\alpha)s^{-\alpha},
\]
and
\[
 D^2\Psi
 =
 q''(d)\nabla d\otimes\nabla d+q'(d)D^2d.
\]
Since $\mathcal A_\lambda$ is independent of $x$,
\begin{align}
 -\operatorname{div}\mathcal A_\lambda(\nabla\Psi)
 &=
 -D\mathcal A_\lambda(\nabla\Psi):D^2\Psi \notag\\
 &\ge
 \nu E(2-\alpha)(1-\alpha)d^{-\alpha}
 -C_1|q'(d)|,
 \label{eq:barrier-operator}
\end{align}
where $C_1$ depends only on $N,L$ and the $C^2$ geometry of
$\partial\Omega$.

We now fix the parameters.  Set
\[
 c_\alpha:=(2-\alpha)(1-\alpha)>0.
\]
First choose $K>0$, depending only on the structural and geometric
constants, so large that
\begin{equation}\label{eq:K-choice-final}
 \nu Kc_\alpha\ge 4C_1+4.
\end{equation}
Next choose $0<\rho\le\rho_0$ sufficiently small that
\begin{equation}\label{eq:rho-choice-final}
 K(2-\alpha)\rho\le\frac12,
 \qquad
 K\rho\le\frac12,
\end{equation}
and, in addition,
\begin{equation}\label{eq:rho-singular-choice-final}
 \frac{\nu Kc_\alpha}{4}\frac{2M}{\rho}\ge H
 \quad\hbox{and}\quad
 \frac{\nu Kc_\alpha}{4}
 \frac{2M}{\rho}\rho^\alpha
 \ge
 \left(\frac{M}{\rho}\right)^{-\alpha}.
\end{equation}
The last two requirements can indeed be achieved by decreasing
$\rho$: the first left-hand side is of order $\rho^{-1}$, whereas in
the second inequality the left-hand side is of order
$\rho^{\alpha-1}$ and the right-hand side is of order $\rho^\alpha$.
Having fixed $K$ and $\rho$, define
\begin{equation}\label{eq:AB-coupled-choice}
 A:=\frac{2M}{\rho},
 \qquad
 E:=KA\rho^\alpha.
\end{equation}
All these constants are independent of $\lambda$.

For $0\le s\le\rho$, we have
\[
 E(2-\alpha)s^{1-\alpha}
 \le
 KA\rho^\alpha(2-\alpha)\rho^{1-\alpha}
 =KA(2-\alpha)\rho
 \le\frac A2.
\]
Consequently,
\begin{equation}\label{eq:qprime-final}
 \frac A2\le q'(s)\le A,
 \qquad 0\le s\le\rho,
\end{equation}
and hence
\begin{equation}\label{eq:Psi-comparable-d}
 \frac A2d\le\Psi\le Ad
 \qquad\hbox{in }\Omega_\rho.
\end{equation}
Moreover, on the inner boundary $\{d=\rho\}$,
\begin{align}
 \Psi
 &=A\rho-E\rho^{2-\alpha}\notag\\
 &=2M(1-K\rho)\ge M.
 \label{eq:inner-boundary-barrier}
\end{align}
Thus $\Psi\ge M$ on the inner boundary $\{d=\rho\}$.

It remains to verify the supersolution inequality.  By
\eqref{eq:barrier-operator}, \eqref{eq:AB-coupled-choice}, and
\eqref{eq:qprime-final},
\begin{align}
 -\operatorname{div}\mathcal A_\lambda(\nabla\Psi)
 &\ge
 \nu KAc_\alpha\rho^\alpha d^{-\alpha}-C_1A\notag\\
 &=
 \nu KAc_\alpha\left(\frac{\rho}{d}\right)^\alpha-C_1A.
 \label{eq:barrier-coupled-lower}
\end{align}
Set $\mathcal T:=\nu KAc_\alpha\rho^\alpha d^{-\alpha}$. Since
$0<d\le\rho$, we have $\mathcal T\ge\nu KAc_\alpha$. Hence, by
\eqref{eq:K-choice-final}, $\frac14\mathcal T\ge C_1A$, and, by the first
condition in \eqref{eq:rho-singular-choice-final},
$\frac14\mathcal T\ge\frac{\nu Kc_\alpha}{4}\frac{2M}{\rho}\ge H$.  On the other hand,
\eqref{eq:Psi-comparable-d} gives
\[
 \Psi^{-\alpha}
 \le
 \left(\frac A2\right)^{-\alpha}d^{-\alpha}.
\]
The second condition in \eqref{eq:rho-singular-choice-final} is exactly
\[
 \frac{\nu KAc_\alpha\rho^\alpha}{4}
 \ge
 \left(\frac A2\right)^{-\alpha},
\]
and therefore
\[
 \frac14\mathcal T=\frac{\nu KAc_\alpha\rho^\alpha}{4}d^{-\alpha}
 \ge \Psi^{-\alpha}.
\]
Combining these estimates, $-\operatorname{div}\mathcal A_\lambda(\nabla\Psi)
\ge\mathcal T-C_1A\ge\frac34\mathcal T$, which yields
\begin{equation}\label{eq:Psi-supersolution}
 -\operatorname{div}\mathcal A_\lambda(\nabla\Psi)
 \ge \Psi^{-\alpha}+H
 \qquad\hbox{in }\Omega_\rho,
\end{equation}
uniformly in $\lambda$.

On $\partial\Omega$ both $w_\lambda$ and $\Psi$ vanish, while on
$\{d=\rho\}$, \eqref{eq:uniform-preboundary} and
\eqref{eq:inner-boundary-barrier} give $w_\lambda\le M\le\Psi$.
Testing the difference between \eqref{eq:abstract-scaled} and
\eqref{eq:Psi-supersolution} with $(w_\lambda-\Psi)^+$ gives
\begin{align*}
 &\int_{\{w_\lambda>\Psi\}}
 \bigl(\mathcal A_\lambda(\nabla w_\lambda)
       -\mathcal A_\lambda(\nabla\Psi)\bigr)
 \cdot\nabla(w_\lambda-\Psi)\,dx\\
 &\qquad\le
 \int_{\{w_\lambda>\Psi\}}
 (w_\lambda^{-\alpha}-\Psi^{-\alpha})
 (w_\lambda-\Psi)\,dx\le0,
\end{align*}
because $s\mapsto s^{-\alpha}$ is decreasing.  The strong
monotonicity implied by \eqref{eq:A-uniform-ellipticity} yields
$w_\lambda\le\Psi$ in $\Omega_\rho$.  Therefore
\[
 w_\lambda\le Ad\qquad\hbox{in }\Omega_\rho.
\]
If $d(x)\ge\rho$, then
\[
 w_\lambda(x)\le M\le\frac M\rho d(x).
\]
Consequently
\[
 w_\lambda(x)\le
 \max\left\{A,\frac M\rho\right\}d(x)
 \qquad\hbox{in }\Omega,
\]
which proves the assertion.
\end{proof}

\begin{lemma}[Uniform singular $C^{1,\beta}$ estimate]
\label{lem:uniform-singular-C1beta}
Under the hypotheses of
Lemma~\ref{lem:uniform-linear-boundary}, there exist
$\beta\in(0,1)$ and $C>0$, independent of $\lambda$, such that
\begin{equation}\label{eq:uniform-C1beta}
 \|w_\lambda\|_{C^{1,\beta}(\overline\Omega)}\le C.
\end{equation}
Consequently,
\begin{equation}\label{eq:u-C1-small}
 \|t_\lambda w_\lambda\|_{C^1(\overline\Omega)}
 \le Ct_\lambda.
\end{equation}
\end{lemma}

\begin{proof}
By Lemma~\ref{lem:uniform-linear-boundary},
\[
 c_0d(x)\le w_\lambda(x)\le C_dd(x).
\]
In particular,
\[
 w_\lambda(x)^{-\alpha}
 \le c_0^{-\alpha}d(x)^{-\alpha}.
\]
Since $0\le h_\lambda\le H$, if
$D=\operatorname{diam}\Omega$, then
\begin{align}
 0\le F_\lambda(x)
 &:=
 w_\lambda(x)^{-\alpha}+h_\lambda(x)\notag\\
 &\le c_0^{-\alpha}d(x)^{-\alpha}+H
 \le C_Fd(x)^{-\alpha},
 \label{eq:F-pure-distance}
\end{align}
where
\[
 C_F:=c_0^{-\alpha}+HD^\alpha
\]
is independent of $\lambda$.

We now verify the structural hypotheses in
\cite[Theorem B.1]{GiacomoniSchindlerTakac}.  From
\[
 \mathcal A_\lambda(\xi)
 =a(t_\lambda^2|\xi|^2)\xi
\]
we have
\begin{equation}\label{eq:DA-lambda-final}
 D\mathcal A_\lambda(\xi)
 =
 a(t_\lambda^2|\xi|^2)I
 +
 2t_\lambda^2a'(t_\lambda^2|\xi|^2)\xi\otimes\xi.
\end{equation}
By hypothesis \eqref{eq:A-uniform-ellipticity} (which, for the
coefficient \eqref{eq7}, holds with $\nu=\nu_0$ and $L=1$ by
Remark~\ref{rem:DA}), uniformly in $\lambda$,
\[
 \nu|\eta|^2
 \le D\mathcal A_\lambda(\xi)\eta\cdot\eta,
 \qquad
 |D\mathcal A_\lambda(\xi)|\le L.
\]
Thus the structural assumptions of
\cite[Theorem B.1]{GiacomoniSchindlerTakac} are satisfied in the
uniformly elliptic case $p=2$.  Moreover,
$\mathcal A_\lambda$ has no explicit $x$-dependence, so the
corresponding continuity assumption in the spatial variable is
automatic.  The datum satisfies
\[
 0\le F_\lambda(x)\le C_Fd(x)^{-\delta},
 \qquad \delta=\alpha\in(0,1),
\]
by \eqref{eq:F-pure-distance}, and the additional boundary hypothesis
required in that theorem,
\[
 0\le w_\lambda(x)\le C_dd(x),
\]
is exactly the upper estimate supplied by
Lemma~\ref{lem:uniform-linear-boundary}.

Therefore \cite[Theorem B.1]{GiacomoniSchindlerTakac} yields an
exponent $\beta\in(0,1)$ and a constant $C>0$, depending only on the
domain, the structural constants $\nu,L$, $\alpha$, $C_F$, and $C_d$
(hence only on $\Omega$, $N$, $\alpha$, $\nu$, $L$, $c_0$, $M$ and $H$), such
that
\[
 \|w_\lambda\|_{C^{1,\beta}(\overline\Omega)}\le C.
\]
Since all these quantities are uniform in $\lambda$, the same
$\beta$ and $C$ work for the whole family.  Finally,
\[
 \|t_\lambda w_\lambda\|_{C^1(\overline\Omega)}
 \le
 t_\lambda\|w_\lambda\|_{C^1(\overline\Omega)}
 \le Ct_\lambda,
\]
as claimed.
\end{proof}

\begin{remark}\label{rem:beta-not-explicit}
For the decontraction argument only a positive uniform H\"older exponent
is needed.  Accordingly, we use precisely the conclusion furnished by
\cite[Theorem B.1]{GiacomoniSchindlerTakac} and do not claim here that
every $\beta<1-\alpha$ is admissible.
\end{remark}

\begin{lemma}[Uniform lower barrier]\label{lem:lower-uniform}
There exists $c_1>0$, independent of $\lambda>0$, such that
\[
u_{0,\lambda}(x)\ge c_1t_\lambda d(x)\qquad\hbox{in }\Omega .
\]
\end{lemma}

\begin{proof}
Let $v_\lambda:=u_{0,\lambda}/t_\lambda$ and
$\mathcal A_\lambda(\xi):=a(t_\lambda^2|\xi|^2)\xi$. Since
$\lambda=t_\lambda^{1+\alpha}$, $v_\lambda$ solves
\[
-\operatorname{div}\mathcal A_\lambda(\nabla v_\lambda)=v_\lambda^{-\alpha}
\quad\hbox{in }\Omega,\qquad v_\lambda=0\quad\hbox{on }\partial\Omega .
\]
Let $\varphi_1\in C^2(\overline\Omega)$ be the positive first Dirichlet
eigenfunction of $-\Delta$, normalized by $\|\varphi_1\|_\infty=1$; by the
Hopf lemma, $\varphi_1\ge k\,d$ for some $k>0$. For $\varepsilon>0$ put
$z:=\varepsilon\varphi_1$. By Remark~\ref{rem:DA},
$\|D\mathcal A_\lambda\|\le1$, hence
\[
-\operatorname{div}\mathcal A_\lambda(\nabla z)
=-D\mathcal A_\lambda(\nabla z):D^2z
\le N\varepsilon\|D^2\varphi_1\|_\infty .
\]
Choose $\varepsilon_0>0$, depending only on $N$, $\alpha$ and $\Omega$,
such that $N\varepsilon_0^{1+\alpha}\|D^2\varphi_1\|_\infty\le1$. Then,
since $z\le\varepsilon_0$,
\[
-\operatorname{div}\mathcal A_\lambda(\nabla z)
\le\varepsilon_0^{-\alpha}\le z^{-\alpha}
\quad\hbox{in }\Omega
\]
for $z=\varepsilon_0\varphi_1$. Testing the difference of the two
problems with $(z-v_\lambda)^+\in H_0^1(\Omega)$ and using that
$s\mapsto s^{-\alpha}$ is decreasing and \eqref{strongmono} (which holds
for $\mathcal A_\lambda$), we get $(z-v_\lambda)^+=0$. Hence
$v_\lambda\ge\varepsilon_0\varphi_1\ge\varepsilon_0k\,d$, and
$c_1:=\varepsilon_0k$ does not depend on $\lambda$.
\end{proof}

\begin{lemma}\label{scaledC1}
There exist $\lambda_0\in(0,\Lambda]$, $\beta\in(0,1)$ and $C_0\ge G$,
independent of $\lambda\in(0,\lambda_0]$, such that the following holds.
Let $R_\lambda:=C_0t_\lambda$ and let $u_{\lambda,R_\lambda}$ be any
solution of \eqref{PRT} with $R=R_\lambda$ supplied by
Proposition~\ref{existtruncated}. Then $u_{\lambda,R_\lambda}$ solves
\eqref{PR},
\[
\|u_{\lambda,R_\lambda}\|_{C^{1,\beta}(\overline\Omega)}\le C_0t_\lambda
\quad\hbox{and, in particular,}\quad
\|u_{\lambda,R_\lambda}\|_{C^1(\overline\Omega)}
\le C_0t_\lambda.
\]
\end{lemma}

\begin{proof}
The constants are fixed in the following order, so that no constant
depends on one chosen later.

\emph{Step 1: structure.} By Remark~\ref{rem:DA}, the family
$\mathcal A_\lambda(\xi):=a(t_\lambda^2|\xi|^2)\xi$ satisfies
\eqref{eq:A-uniform-ellipticity} with $\nu=\nu_0$ and $L=1$, for every
$\lambda$.

\emph{Step 2: barriers.} Put $c_2:=B_0$ and let $c_1$ be given by
Lemma~\ref{lem:lower-uniform}. Both depend only on $\Omega$, $N$,
$\alpha$ (and on the fixed point $x_0$).

\emph{Step 3: the a priori constant.} Apply
Lemmas~\ref{lem:uniform-linear-boundary} and
\ref{lem:uniform-singular-C1beta} with $\nu=\nu_0$, $L=1$, $c_0=c_1$,
$M=c_2$ and $H=1$. They give $\beta\in(0,1)$ and $C_\sharp>0$, depending
only on these quantities, on $\alpha$ and on $\Omega$, such that every
positive $w\in H_0^1(\Omega)\cap L^\infty(\Omega)$ satisfying
\eqref{eq:abstract-scaled} with
\begin{equation}\label{eq:hyp-scaled}
0\le h_\lambda\le1,\qquad c_1d\le w\le c_2,
\end{equation}
obeys $\|w\|_{C^{1,\beta}(\overline\Omega)}\le C_\sharp$. Define
\[
C_0:=\max\{C_\sharp,G\}.
\]

\emph{Step 4: smallness of $\lambda$.} Since $r_1,r_2>1$, we can choose
$\lambda_0\in(0,\Lambda]$ such that, for $0<\lambda\le\lambda_0$,
\begin{equation}\label{eq:lambda0}
k_1c_2^{r_1}t_\lambda^{r_1-1}+k_2C_0^{r_2}t_\lambda^{r_2-1}\le1
\qquad\hbox{and}\qquad C_0t_\lambda<1 .
\end{equation}

\emph{Step 5: conclusion.} Let $0<\lambda\le\lambda_0$ and
$u=u_{\lambda,R_\lambda}$. Since $R_\lambda=C_0t_\lambda\ge Gt_\lambda$,
Proposition~\ref{location} applies: $u$ solves \eqref{PR} and
$\underline u_\lambda\le u\le\overline u_\lambda$. Write $u=t_\lambda w$.
Since $\lambda=t_\lambda^{1+\alpha}$, equation \eqref{PR} becomes
\begin{equation}\label{scaled}
-\operatorname{div}\left(a(t_\lambda^2|\nabla w|^2)\nabla w\right)
=w^{-\alpha}
+h_\lambda,\qquad
h_\lambda:=t_\lambda^{-1}f_{R_\lambda}
(x,t_\lambda w,t_\lambda\nabla w).
\end{equation}
By Lemma~\ref{lem:lower-uniform} and
$\overline u_\lambda=t_\lambda\vartheta\le t_\lambda B_0$, we have
$c_1d\le w\le c_2$. Moreover, since $|\pi_{R_\lambda}(\xi)|\le C_0t_\lambda$,
assumption \eqref{f} and \eqref{eq:lambda0} give
\[
0\le h_\lambda
\le t_\lambda^{-1}\bigl(k_1(t_\lambda c_2)^{r_1}+k_2(C_0t_\lambda)^{r_2}\bigr)
=k_1c_2^{r_1}t_\lambda^{r_1-1}+k_2C_0^{r_2}t_\lambda^{r_2-1}\le1 .
\]
Thus \eqref{eq:hyp-scaled} holds and Step~3 yields
$\|w\|_{C^{1,\beta}(\overline\Omega)}\le C_\sharp\le C_0$. Since
$u=t_\lambda w$,
\[
\|u\|_{C^{1,\beta}(\overline\Omega)}\le C_\sharp t_\lambda\le C_0t_\lambda,
\]
as claimed.
\end{proof}

\begin{proposition}\label{decontract}
For $0<\lambda\le\lambda_0$, every solution
$u_{\lambda,R_\lambda}$ as in Lemma~\ref{scaledC1} satisfies
\[
|\nabla u_{\lambda,R_\lambda}(x)|\le R_\lambda
\qquad\hbox{for every }x\in\overline\Omega.
\]
Hence
\[
f_{R_\lambda}(x,u_{\lambda,R_\lambda},
\nabla u_{\lambda,R_\lambda})
=
f(x,u_{\lambda,R_\lambda},\nabla u_{\lambda,R_\lambda}),
\]
and $u_{\lambda,R_\lambda}$ solves the uncontracted auxiliary problem
\eqref{eq6}.
\end{proposition}

\begin{proof}
This is immediate from Lemma \ref{scaledC1} and the definition
$R_\lambda=C_0t_\lambda$.
\end{proof}

\section{Proof of Theorem \ref{main result}}

\begin{proof}

Let $\lambda^*:=\lambda_0$, where $\lambda_0\le\Lambda$ is given by
Lemma~\ref{scaledC1}; in particular $C_0t_\lambda<1$ for
$0<\lambda\le\lambda^*$, by \eqref{eq:lambda0}.
Let $u_\lambda=u_{\lambda,R_\lambda}$, which exists by
Proposition~\ref{existtruncated}.  By Proposition
\ref{decontract}, $u_\lambda$ solves \eqref{eq6}.  Since
$\|\nabla u_\lambda\|_\infty<1$, the definition \eqref{eq7} gives
\[
a(|\nabla u_\lambda|^2)
=\frac{1}{\sqrt{1+|\nabla u_\lambda|^2}}
\qquad\hbox{in }\overline\Omega.
\]
Therefore \eqref{eq6} is exactly the original prescribed mean-curvature
problem \eqref{P1}.  Positivity follows from
$u_\lambda\ge\underline u_\lambda>0$ in $\Omega$, and Lemma
\ref{scaledC1} yields
\[
\|u_\lambda\|_{C^1(\overline\Omega)}
\le C_0\lambda^{1/(1+\alpha)}\longrightarrow0
\qquad\hbox{as }\lambda\to0^+.
\]
Finally, Lemma~\ref{scaledC1} gives
$u_\lambda\in C^{1,\beta}(\overline\Omega)$ with $\beta\in(0,1)$
independent of $\lambda$.  This proves Theorem
\ref{main result}.
\end{proof}


\noindent{\sc Acknowledgment}. A. L. A. de Araujo has been partially supported by FAPEMIG and CNPq.

\vspace{0.5cm}

\section*{Declarations}
\noindent{\bf Conflict of interests}. The authors have no competing interests to declare that are relevant to the content of
this article.

\begin{thebibliography}{10}


\bibitem{AFMsys2024}
A. de Araujo, L. F. O. Faria and D. Motreanu,
{\em Quasilinear systems with unbounded variable exponents and convection terms},
Bull. Braz. Math. Soc. (N.S.) \textbf{55} (2024), Article 17.
doi:10.1007/s00574-024-00391-x.

\bibitem{BOP2024}
F. Balducci, F. Oliva and F. Petitta,
{\em Finite energy solutions for nonlinear elliptic equations with competing
gradient, singular and $L^1$ terms},
J. Differential Equations \textbf{391} (2024), 334--369.
doi:10.1016/j.jde.2024.02.002.

\bibitem{FZ2024}
R. Filippucci and Y. Zheng,
{\em Existence and nonexistence of solutions for the mean curvature equation
with weights},
Math. Ann. \textbf{390} (2024), 5821--5861.
doi:10.1007/s00208-024-02900-1.

\bibitem{GOP2024}
D. Giachetti, F. Oliva and F. Petitta,
{\em Bounded solutions for non-parametric mean curvature problems with
nonlinear terms},
J. Geom. Anal. \textbf{34} (2024), Article 265.
doi:10.1007/s12220-024-01715-5.

\bibitem{OP2025}
F. Oliva and F. Petitta,
{\em Singular elliptic PDEs: an extensive overview},
Partial Differential Equations Appl. \textbf{6} (2025), Article 6.
doi:10.1007/s42985-024-00308-9.

\bibitem{OPS2024}
F. Oliva, F. Petitta and S. Segura de Le\'on,
{\em The role of absorption terms in Dirichlet problems for the prescribed
mean curvature equation},
Nonlinear Differential Equations Appl. \textbf{31} (2024), Article 53.
doi:10.1007/s00030-024-00936-5.


\bibitem{AFM2024}
A. de Araujo, L. Faria and D. Motreanu,
{\em Positive solutions of nonlinear elliptic equations involving
unbounded variable exponents and convection term},
Partial Differential Equations and Applications \textbf{5} (2024), 27.


\bibitem{CGMS} Ph. Cl\'ement, M. Garc\'ia-Huidobro, R. Man\'asevich and K. Schmitt, {\em Mountain pass type solutions for quasilinear elliptic equations}. Calc. Var. Partial Differential Equations \textbf{11} (2000) 33--62.

\bibitem{CoffmanZ} C. V. Coffman and W. K. Ziemer, {\em A prescribed mean curvature problem on domains without radial symmetry}, SIAM J. Math. Anal. \textbf{22} (1991) 982--990.

\bibitem{FR} G. M. Figueiredo and V. D. R\u{a}dulescu, {\em Positive solutions of the prescribed mean curvature equation with exponential critical growth}. Ann. Mat. Pura Appl. (4) \textbf{200} (2021) 2213--2233. doi:10.1007/s10231-021-01077-7.

\bibitem{Finn} R. Finn, Equilibrium capillary surfaces, Grundlehren Math. Wiss. Band \textbf{284}. Springer-Verlag (1986).

\bibitem{GCS} J. V. Gon\c{c}alves, M. L. Carvalho and C. A. Santos, {\em Quasilinear elliptic systems convex-concave singular terms and $\Phi$--Laplacian operator}. Differential and Integral Equations \textbf{31} (2018) 231--256.


\bibitem{Giusti} E. Giusti, Minimal surfaces and functions of bounded variation. Monographs in Mathematics \textbf{80}. Birkh\"auser Verlag (1984).


\bibitem{HabetsO} P. Habets and P. Omari, {\em Positive solutions of an indefinite prescribed mean curvature problem on a general domain}. Adv. Nonlinear Stud. \textbf{4} (2004) 1--13.
	
\bibitem{KOF} A. Kufner, O. John and S. Fu\v{c}\'{\i}k, Function Spaces, Noordhoff International Publishing, Leyden; Academia, Prague (1977).

\bibitem{Le} V. K. Le, {\em Variational method based on finite dimensional approximation in a generalized prescribed mean curvature problem}. J. Diff. Equations \textbf{246} (2009) 3559--3578.

	
\bibitem{GiacomoniSchindlerTakac}
J. Giacomoni, I. Schindler and P. Tak\'a\v{c},
{\em Sobolev versus H\"older local minimizers and existence of
multiple solutions for a singular quasilinear equation},
Ann. Sc. Norm. Super. Pisa Cl. Sci. (5) \textbf{6} (2007), 117--158.

\bibitem{delPino}
M. A. del Pino,
{\em A global estimate for the gradient in a singular elliptic boundary
value problem},
Proc. Roy. Soc. Edinburgh Sect. A \textbf{122} (1992), 341--352.
doi:10.1017/S0308210500021144.

\bibitem{GuiLin}
C. Gui and F. H. Lin,
{\em Regularity of an elliptic problem with a singular nonlinearity},
Proc. Roy. Soc. Edinburgh Sect. A \textbf{123} (1993), 1021--1029.
doi:10.1017/S030821050002970X.

\bibitem{Lie} G. M. Lieberman, {\em Boundary regularity for solutions of degenerate elliptic equations}. Nonlinear Anal. \textbf{12} (1988) 1203--1219.

\bibitem{Lopez}  R. L\'opez, {\em An existence theorem of constant mean curvature graphs in Euclidean space}. Glasgow Math. J. \textbf{44} (2002) 455--461.

\bibitem{Lorca1}  S. Lorca and M. Montenegro,  {\em Multiple solutions for the mean curvature equation}. Topol. Methods Nonlinear Anal. \textbf{35} (2010) 61--68.

\bibitem{Lorca2}  S. Lorca and M. Montenegro,  {\em Height estimates for graph parametrized surfaces}. Bull. Sci. Math. \textbf{136} (2012) 848--856.

\bibitem{Mont.} M. Montenegro, {\em Strong maximum principles for supersolutions of quasilinear elliptic equations}. Nonlinear Anal. \textbf{37} (1999) 431--448.


\bibitem{ObersnelO} F. Obersnel and P. Omari, {\em Positive solutions of the Dirichlet problem for the prescribed mean curvature equation}. J. Diff. Equations \textbf{249} (2010) 1674--1725.

\bibitem{RR} M. N. Rao and Z. D. Ren, Theory of Orlicz Spaces, Marcel Dekker, New York (1985).

\bibitem{Ripoll} J. Ripoll,  {\em Some existence results and gradient estimates of solutions of the Dirichlet problem for the constant mean curvature equation in convex domains}. J. Differential
Equations \textbf{181} (2002) 230--241.

\bibitem{Serrin}  J. Serrin,  {\em The problem of Dirichlet for quasilinear elliptic differential equations with many independent variables}. Philos. Trans. Roy. Soc. London Ser. A \textbf{264} (1969) 413--496.

\bibitem{Tersenov} A. S. Tersenov,  {\em On quasilinear non-uniformly elliptic equations in some non-convex domains}. Comm. Part. Diff. Equations \textbf{23} (1998) 2165--2185.


	
\end{thebibliography}

\end{document}

